% ---------------------------------------------------------------------------
% Shared preamble for the standalone code listings in this directory.
%
% The font sizes and the widths below mirror those set by nowfnt.cls /
% main.tex so that the compiled PDFs can be dropped into the main document at
% 1:1 scale.  See README.md.
% ---------------------------------------------------------------------------

\usepackage[T1]{fontenc}
\usepackage{lmodern}
\usepackage{listings}
\usepackage{xcolor}

% Font size overrides copied from nowfnt.cls so that \scriptsize (basicstyle)
% and \tiny (line numbers) match the main document exactly.
\makeatletter
\renewcommand\normalsize{\@setfontsize\normalsize\@xipt\@xivpt}
\renewcommand\scriptsize{\@setfontsize\scriptsize\@viipt\@viiipt}
\renewcommand\tiny{\@setfontsize\tiny\@vpt\@vipt}
\makeatother
\normalsize

% --- Geometry --------------------------------------------------------------
% A listing occupies three horizontal bands, only the middle of which is the
% code itself:
%
%   |<- NumGutter ->|<---------- BodyWidth ---------->|<- Over ->|
%   |   1 2 3 ...   |  // the code, frame=single      |          |
%           |<-Over->|                                |
%
% The line numbers hang to the LEFT of the frame, and the frame rule is drawn
% \lstFrameOver outside the body on BOTH sides.  Anything outside the
% standalone bounding box is cropped away, so the box has to span all three
% bands -- that is what \lstTotalWidth is for.

% \textwidth of the ROB journal layout.  The code body keeps exactly this
% measure, so line wrapping matches the old inline lstlisting output.
\newlength{\lstBodyWidth}
\setlength{\lstBodyWidth}{335.74251pt}

% Room for the line numbers, which hang to the left of the frame.
\newlength{\lstNumGutter}
\setlength{\lstNumGutter}{20pt}

% How far frame=single reaches outside the body: framesep + framerule.
\newlength{\lstFrameOver}
\setlength{\lstFrameOver}{3.4pt}

\newlength{\lstTotalWidth}
\setlength{\lstTotalWidth}{\lstBodyWidth}
\addtolength{\lstTotalWidth}{\lstNumGutter}
\addtolength{\lstTotalWidth}{\lstFrameOver}

% Every listing goes inside this, which is exactly \lstTotalWidth wide, so that
% standalone crops the page to the whole listing rather than to the code body
% alone.  (Neither the standalone `varwidth' option nor a bare minipage is
% enough: both end up measuring the width of the code body, leaving the
% right-hand frame rule off the page.  The empty, zero-height \hbox pins the
% box to its full width without adding any vertical space.)
\newenvironment{listingbox}
  {\begin{minipage}{\lstTotalWidth}}
  {\par\nointerlineskip\hbox to \lstTotalWidth{}%
   \end{minipage}}

% c++ listing format (kept identical to the definition in main.tex, apart from
% the two margins, which position the code within \lstTotalWidth)
\lstdefinestyle{C++Style}{
  language=C++,
  basicstyle=\scriptsize\ttfamily, % Smaller monospaced font
  keywordstyle=\color{blue}, % Keywords in blue
  commentstyle=\color{green!50!black}, % Comments in dark green
  stringstyle=\color{red}, % Strings in red
  showstringspaces=false, % Don't show spaces in strings
  numbers=left, % Line numbers on the left
  numberstyle=\tiny\color{gray}, % Style for line numbers
  frame=single, % Add a frame around the listing
  breaklines=true, % Allow line breaking
  backgroundcolor=\color{black!5}, % Light gray background
  upquote=true,  % Don't convert ' into curly smart quotes
  xleftmargin=\lstNumGutter,   % leave room for the line numbers
  xrightmargin=\lstFrameOver,  % leave room for the right frame rule
}
